IMPORTANT EXERCISE (THE FHHF THEOREM). This result can take you far, and perhaps for that reason it has sometimes been called the Fernbahnhof (FernbaHnHoF) Theorem, notably in [Vak1, Exer. 1.5.I]. Suppose $F : \mathcal{A} \to \mathcal{B}$ is a covariant functor of abelian categories, and $C_\bullet$ is a complex in $\mathcal{A}$. (a) ($F$ right-exact yields $F H_i(C_\bullet) \to H_i(F C_\bullet)$.) If $F$ is right-exact, describe a natural morphism $F H_i(C_\bullet) \to H_i(F C_\bullet)$. (More precisely, for each $i$, the left side is $F$ applied to the cohomology at piece $i$ of $C_\bullet$, while the right side is the cohomology at piece $i$ of $F C_\bullet$.) (b) ($F$ left-exact yields $H_i(F C_\bullet) \to F H_i(C_\bullet)$.) If $F$ is left-exact, describe a natural morphism $H_i(F C_\bullet) \to F H_i(C_\bullet)$. (c) ($F$ exact yields $F H_i(C_\bullet) \cong H_i(F C_\bullet)$.) If $F$ is exact, show that the morphisms of (a) and (b) are inverses and thus isomorphisms. Hint for (a): Use the natural isomorphism $F(\operatorname{coker} d_{i+1}) \cong \operatorname{coker}(F d_{i+1})$. Then use the first line of (1.5.6.4) to give an epimorphism $F(\operatorname{im} d_i) \twoheadrightarrow \operatorname{im}(F d_i)$. Then use the second line of (1.5.6.4) to give the desired map $F H_i(C_\bullet) \to H_i(F C_\bullet)$. While you are at it, you may as well describe a map for the fourth member of the quartet $\{\operatorname{coker},\operatorname{im},H,\ker\}$: $F(\ker d_i) \to \ker(F d_i)$.
